\problemname{Tina poängtavla}

ICPC är en programmeringstävling där man får fem timmar på sig att lösa så många problem som möjligt.
I denna får man inte poäng per inskick, utan antingen får man \textit{Accepted} eller inte.

I slutet av tävlingen etablerar man en rangordning av lag på följande vis.
Först räknar man hur många problem som varje lag löst, låt detta betecknas av $p$.
Sedan beräknar man varje lags \textit{inskick-straff}.
För varje problem där ett lag får Accepted är inskick-straffet lika med antal felaktiga inskick laget gjorde på problemet.
Inskick-straffet för ett lag är summan av inskick-straff över alla problem de löste, vilket vi betecknar med $t$.
Varje lags resultat kan därmed beskrivas med talparet $(p,t)$.
Notera att inskick på problem som ett lag inte löser under hela tävlingen inte ökar inskick-straffet.

Rangen av ett visst lag är antalet lag som är \textit{strikt bättre}.
Ett lag $i$ är strikt bättre över ett annat lag $j$ om antingen
$$(p_i>p_j) \text{ eller } (p_i=p_j \text{ och } t_i < t_j)$$
håller.

I nedanstående tabell följer ett exempel på rangen olika lag får, givet en uppsättning $p$- och $t$-värden.

\begin{center}
\noindent
\begin{tabular}{| c | c | c |}
  \hline
  Rang & $p$ & $t$ \\ \hline
  $0$ & $10$ & $500$ \\ \hline
  $1$ & $9$ & $1000$ \\ \hline
  $2$ & $8$ & $500$ \\ \hline
  $3$ & $7$ & $300$ \\ \hline
  $3$ & $7$ & $300$ \\ \hline
  $5$ & $6$ & $250$ \\ \hline
\end{tabular}
\end{center}

Notera alltså att om två lag har exakt samma resultat får de samma rang.

De första fyra timmarna under tävlingen får varje lag se vare sig deras inskick fick Accepted.
De kan även se resultaten av andra lags inskick genom poängtavlan.

Den sista timmen blir poängtavlan \textit{fryst}, och därmed kan lagen endast se resultatet av sina egna inskick.
Varje lag kan fortfarande se när nya inskick sker av andra lag, men inte deras resultat.

När tävlingen är slut, är det dags att \textit{tina} poängtavlan.
Detta sker genom att resultatet av varje inskick avslöjas en åt gången.
Beroende på ordningen detta sker kan det bli mer eller mindre spännande för lagen.
Detta beror på att man i första hand etablerar rangen av ett lag baserat på hur många problem de har löst.
Exempelvis, om ett visst lag har löst $4$ problem och det avslöjas att alla andra i tävlingen har löst fler,
så tappar laget intresset.
Om inskicken istället hade avslöjats i en ordning så att det tar längre tid för de att få reda på detta blir det självfallet mer spännande.

Lidia har som målsättning att göra ICPC mer spännande.
Hon har därför föreslagit en ordning på hur resultat av inskick ska avslöjas.
Givet denna ordningen vill hon för varje lag beräkna antalet inskick som avslöjas innan de tappar intresset.

Mer exakt, anta att varje lag efter varje inskickning betraktar alla möjliga utfall av att poängtavlan tinas färdigt.
Om ett visst lag har samma slutposition i alla dessa tappar det laget intresset.
Lidia vill då veta för varje lag, antalet inskickningar som avslöjas innan de tappade intresset.

För att ordningen ska vara rimlig garanterar hon att varje lags inskick på ett givet problem är ordnade på tid,
så att tidigare inskick kommer först.

\section*{Indata}
Den första raden innehåller tre heltal $N$, $P$, $H$ och $F$
($1 \leq N \leq 2 \cdot 10^5$, $1 \leq P \leq 15$, $0 \leq H \leq 2 \cdot 10^5$, $1 \leq F \leq 2 \cdot 10^5$),
antalet lag, antalet problem, antalet inskick innan frysen och antalet inskick under frysen.

De följande $H$ raderna beskriver inskicken som skedde före frysen.
Varje sådant inskick beskrivs av heltalen $i$, $p$ och $v$
($1 \leq i \leq N$, $1 \leq p \leq P$, $v \in \{\texttt{``AC''} , \texttt{``WA''}\}$),
som beskriver att laget med index $i$ skickade in till problem $p$,
vars resultat var $v$. Om $v$ är \texttt{``AC''} betyder det att inskicket fick Accepted, annars fick det inte det.
Dessa inskick är ordnade i kronologisk ordning, det vill säga är de sorterade på tid.

De följande $F$ raderna beskriver inskicken som skedde före frysen.
Formatet på dessa är identiska till inskicken som skedde innan poängtavlan frystes.
Dessa är givna i ordningen som Lidia har föreslagit.
Det är garanterat att problemen är ordnade så att varje lags inskick på ett visst problem kommer i kronologisk ordning.
Hon ger inga garantier utöver detta.

\textbf{Not:} det är dessutom garanterat att ett lag aldrig kommer göra inskick till ett problem de redan löst,
varken innan eller under det att poängtavlan är fryst.

\section*{Utdata}
För varje lag, skriv ut hur många submissions det tog för de att tappa intresset.
Skriv ut dessa värden mellanslagsseparerade på en rad.
Notera att om ett lag vet sin slutposition innan tiningen påbörjats ska du skriva ut $0$ för det laget.

\section*{Poängsättning}
Din lösning kommer att testas på en mängd testfallsgrupper.
För att få poäng för en grupp så måste du klara alla testfall i gruppen.

\noindent
\begin{tabular}{| l | l | p{12cm} |}
  \hline
  \textbf{Grupp} & \textbf{Poäng} & \textbf{Gränser} \\ \hline
  $1$    & $7$        & $N,F \leq 50$, $H=0$, $P=1$ och varje lag gör som mest ett inskick per problem. \\ \hline
  $2$    & $9$        & $N,F \leq 50$, $P=1$ och varje lag gör som mest ett inskick per problem. \\ \hline
  $3$    & $15$       & $N,F,H \leq 50$ och varje lag gör som mest ett inskick per problem. \\ \hline
  $4$    & $30$       & $N,F,H \leq 50$ \\ \hline
  $4$    & $80$       & Inga ytterligare begränsningar. \\ \hline
\end{tabular}

% Add description of sample cases
\section*{Förklaring av exempelfall}
Exempelfall 1: innan poängtavlan blir fryst vet lag $2$ redan att de kommer sluta på $1$ problem löst med $0$ inskick-straff.
Därmed vet de redan sin slutgiltiga rang.
Lag $1$ vet inte om lag lag $2$ kommer lösa problemet eller inte, men även om det andra laget löser problemet,
så kan deras resultat inte försämra rangen av lag $1$.

Exempelfall 2: likt exempelfall $1$, så har lag $2$ redan all information.
Lag $1$ kan däremot inte vara säkra på sin slutgiltiga plats direkt efter tävlingen, eftersom om lag $2$ löste problemet så tappar lag $1$ en rang.
När resultatet väl klagjorts blir de säkra.

Exempefall 3: även om det inte verkar så kan lag $1$ lista ut direkt att de har vunnit tävlingen.
Eftersom ett lag aldrig skickar in på problem de redan löst vet lag $1$ att det enda inskicket från som lag 2 som kan få Accepted är den sista.
Även om de får Accepted på denna, så kommer det inte påverka deras rang.
